\documentclass[11pt,a4paper]{article}

\usepackage[utf8]{inputenc}
\usepackage{amsmath}
\usepackage{graphicx}
\usepackage[colorlinks=true,linkcolor=blue,citecolor=blue,urlcolor=blue]{hyperref}
\usepackage{natbib}
\usepackage[margin=1in]{geometry}

\graphicspath{{figures/}}

\title{ {{TITLE}} }
\author{ {{AUTHORS}} }
\date{ {{DATE}} }

\begin{document}
\maketitle

\begin{abstract}
% TODO_ABSTRACT: write a 4-6 sentence abstract.
% Stance hint from skill input: {{STANCE_HINT}}
% Sources: outputs/repro-*/report.md (if reproduction was performed),
%          outputs/slr-*.md (if literature review was performed).
\end{abstract}

\section{Introduction}
% TODO_INTRO: motivate the problem; cite key references via \citep{key}
%             (key must exist in references.bib). If references.bib is empty,
%             write descriptive prose without citations and add inline
%             "% needs citation" markers.

\section{Related Work}
% TODO_RELATED: synthesize from outputs/slr-*.md.
%               Use \citep{key} for inline cites; key must exist in references.bib.
%               If references.bib is empty, prose only with "% needs citation".

\section{Method}
% TODO_METHOD: describe approach;
%              source: outputs/repro-*/report.md (Method section).
%              Include any scale-down decisions explicitly.

\section{Experiments}
% TODO_EXPERIMENTS: present quantitative results.
%                   Sources: outputs/repro-*/comparison.json (paper_value, our_value, delta)
%                            workspace/analysis-*/*.csv (if data-analysis was performed).
%                   Use the figure-include lines below; uncomment ones you cite.

% Figure includes (uncomment as you write Experiments prose):
{{FIGURE_INCLUDES}}

\section{Conclusion}
% TODO_CONCLUSION: 1-paragraph summary of contribution + limitations + future work.

{{BIBLIOGRAPHY_LINE}}

\end{document}
